\subsection{数学速查：遇到符号卡住时先看这里}
\phantomsection\label{app:math-tools}
本书建议从高中代数出发；后面的梯度和 Attention 还会使用
高中三角函数、指数与对数，以及一点初等微积分。
五题自检帮助你选择起点：先补齐薄弱的运算，再回到正文算例。

\begin{center}
\small
\renewcommand{\arraystretch}{1.36}
\tagpdfsetup{table/header-rows={1}}
\begin{tabularx}{\linewidth}{@{}p{31mm}X@{}}
\toprule
符号 & 最简单的意思与算例\\
\midrule
$\sum_{i=1}^3 a_i$ & 三项相加：$a_1+a_2+a_3$\\
$\exp(t)=e^t$ & 以自然常数 $e\approx2.71828$ 为底的指数；$e^0=1$\\
$\ln x$ & 自然对数，以 $e$ 为底；$\ln(e)=1$，$\ln1=0$\\
$f'(x),\ \frac{\mathrm df}{\mathrm dx}$ & 小幅改变 $x$ 时，函数输出变化的局部比率\\
$\frac{\partial L}{\partial w}$ & 多个变量中只改变 $w$ 时的局部变化率\\
$\nabla_{\!\theta}L$ & 各参数偏导构成的向量或张量集合\\
$\cos\theta$ & 两条向量夹角的余弦；余弦定理属高中三角知识\\
$\R^{T\times d}$ & $T$ 行 $d$ 列矩阵，本书主线每一行代表一个 Token\\
$\top$ & 转置：行与列交换，$(AB)^\top=B^\top A^\top$\\
$\odot$ & 相同位置的元素逐一相乘，不是矩阵乘法\\
$\approx$ & 数值近似而非精确相等，例如 $1/3\approx0.333$\\
\bottomrule
\end{tabularx}
\end{center}

\subsubsection*{五题入门自检：能直接开读吗？}
\phantomsection\label{app:math-entry}
花五分钟独立完成，不用记新术语。五题都能说明理由，就可以直接开读；
卡在哪项，先复习上表对应的运算。第 5 题只要求一元导数，链式法则会在正文继续练习。
\begin{enumerate}
\item 解 $2w+1=5$。
\item 计算 $(-2)^2+1$，并说明它与 $-2^2+1$ 的区别。
\item 对 $f(x)=2x+1$，求 $f(-1),f(0),f(1)$。
\item 计算 $\ln(e^2)$，并求满足 $e^u=1$ 的 $u$。
\item 对 $g(w)=(w-2)^2$，求 $g'(w)$ 与 $g'(1)$；应增大还是减小 $w$ 才先靠近最低点？
\end{enumerate}
\hyperref[answer:math-entry]{核对结果与补读建议：附录 D，第~\pageref*{answer:math-entry}~页。}

\subsubsection*{导数、链式法则与一次更新}

如果 $f(x)=x^2$，导数为 $f'(x)=2x$；
例如 $x=3$ 时斜率为 $6$。
如果 $L=(2w-5)^2/2$，令 $u=2w-5$，则
\[
\frac{\mathrm dL}{\mathrm dw}
=\frac{\mathrm dL}{\mathrm du}\frac{\mathrm du}{\mathrm dw}
=u\cdot2=2(2w-5).
\]
这就是\textbf{链式法则}：沿计算路径把局部导数相乘。
学习率 $\eta$ 的梯度下降一步是
$w_{\rm new}=w-\eta\,\mathrm dL/\mathrm dw$。

\subsubsection*{一张对照表解决 Token ID 和矩阵下标}

教学词表 ID、位置与数学行列都从 $1$ 起，直接按编号查表：
\[
\begin{array}{c|ccc}
\text{词表 ID／数学列号} & 1 & 4 & 5\\\hline
\text{Python 列下标} & 0 & 3 & 4
\end{array}
\]
“花”（ID 4）的第一位置概率是 $\Pi_{1,4}$。
写 Python 时，在数组索引处统一减 $1$，读取 \texttt{probs[0,3]}。
脚本保留从 $1$ 起的 ID，并将 \texttt{ids - 1} 用于查表、
\texttt{targets - 1} 用于交叉熵标签。

\subsubsection*{不要把矩阵的行列和向量方向混用}

例如 $\vect x^\top=(1,2)$ 作为\textbf{行向量}，
$W=\begin{bmatrix}1&0\\0&3\end{bmatrix}$，
则 $\vect x^\top W=(1,6)$；如果把
$\vect x=(1,2)^\top$ 改写成列向量，同一映射写为
$W^\top\vect x=(1,6)^\top$。
权重矩阵也要转置，不能只把输入横过来。



\begin{conclusion}
先认清：每个符号是什么、尺寸能否相乘、数字能否代入核对。
到梯度和 Softmax 时，你无需预先记住一整套高等数学；
每遇到一个新运算，都回到这里的小例子。
\end{conclusion}
